%Chargement des paquets
\usepackage{amsmath}
\usepackage{amsfonts}
\usepackage{amssymb}
\usepackage{enumerate}
\usepackage{mathtools}
\usepackage{bbm}
\usepackage{xparse, etoolbox}
\usepackage{enumerate}
\usepackage{enumitem}
\usepackage{mathabx}
\usepackage{babel}[french]

%environment
\usepackage{amsthm}
\usepackage{keytheorems}
\usepackage{hyperref} 

%Interface théorème
\newkeytheorem{lem}[
	name = Lemme
]

\newkeytheorem{prop}[
	name = Proposition
]

\newkeytheorem{thm}[
	name = Théorème
]

\newkeytheorem{coro}[
	name = Corollaire
]

\renewcommand*{\proofname}{Démonstration}

\theoremstyle{definition}
\newtheorem{defn}{Définition}

\theoremstyle{definition}
\newtheorem*{nota}{Notation}

\theoremstyle{definition}
\newtheorem*{limi}{Liminaire}

\theoremstyle{remark}
\newtheorem{rmq}{Remarque}

\theoremstyle{plain}
\newtheorem*{fait}{Fait}


%Ensembles de nombres
\newcommand{\N} {\mathbb{N}}
\newcommand{\Ne}{\N^\ast}
\newcommand{\Z} {\mathbb{Z}}
\newcommand{\Q} {\mathbb{Q}}
\newcommand{\R} {\mathbb{R}}
\newcommand{\Rb}{\overline{\mathbb{R}}}
\newcommand{\Rp}{\R_+}
\newcommand{\Rm}{\R_-}
\newcommand{\K} {\mathbb{K}}
\newcommand{\Ket} {\mathbb{K^{\ast}}}
\newcommand{\Cx}{\mathbb{C}}

%Ensembles, tribus
\newcommand{\A}{\mathbb{A}}
\renewcommand{\P}{\mathcal{P}}
\newcommand{\T}{\mathcal{T}}
\newcommand{\F}{\mathcal{F}}
\newcommand{\C} {\mathcal{C}}
\newcommand{\ouv}{\mathcal{O}}
\newcommand{\B}{\mathcal{B}}
\newcommand{\M}{\mathcal{M}}
\newcommand{\etag}{\mathcal{E}}
\newcommand{\etagOp}{\etag^{+}(\Omega)}
\newcommand{\etagO}{\etag(\Omega)}
%%Lp avant quotientage
\newcommand{\Lic}{\mathcal{L}^1}
\newcommand{\Lpc}{\mathcal{L}^p}
\newcommand{\Linfc}{\mathcal{L}^{\infty}}
\newcommand{\Lifull}{\Lic(\Omega, \B, m)}
\newcommand{\Lpfull}{\Lpc(\Omega, \B, m)}
\newcommand{\Linffull}{\Linfc(\Omega, \B, m)}
%%Lp après quotientage
\newcommand{\Li}{L^1}
\newcommand{\Ld}{L^2}
\newcommand{\Lp}{L^p}
\newcommand{\Linf}{L^{\infty}}
\newcommand{\Listd}{\Li(\Omega, m)} 
\newcommand{\Ldstd}{\Ld(\Omega, m)} 
\newcommand{\Lpstd}{\Lp(\Omega, m)} 

%Quelques opérateurs
\DeclareMathOperator{\codim}{codim}
\DeclareMathOperator{\rg}{rg}
\DeclareMathOperator{\tr}{tr}
\DeclareMathOperator{\vect}{Vect}
\DeclareMathOperator{\ima}{Im}
\DeclareMathOperator{\id}{Id}
\DeclareMathOperator{\sign}{sign}
\DeclareMathOperator{\ext}{ext}
\newcommand{\ind}{\mathbbm{1}}
\newcommand*{\spr}[2]{\left(\,#1\,\middle|\,#2\,\right)}
\newcommand{\interior}[1]{%
  {\kern0pt#1}^{\mathrm{o}}%
}
\DeclareMathOperator{\card}{card}
\DeclareMathOperator{\ppcm}{ppcm}

%Commandes de pure flemme
\newcommand{\veps}{\varepsilon}
\newcommand{\intab}{\int_{a}^{b}}
\newcommand{\intR}{\int_{-\infty}^{\infty}}
\newcommand{\intinf}{\int_{- \infty}^{+ \infty}}
\newcommand{\equi}{\Leftrightarrow}
\newcommand{\Kev}{$\K$-espace vectoriel }
\newcommand{\Kevs}{$\K$-espaces vectoriels }
\newcommand{\Rev}{$\R$-espace vectoriel }
\newcommand{\Revs}{$\R$-espaces vectoriels }
\newcommand{\Cev}{$\Cx$-espace vectoriel }
\newcommand{\Cevs}{$\Cx$-espaces vectoriels }
\newcommand{\evn}{(E, \norm{.})}
\newcommand{\interf}[2]{[#1,#2]}
\newcommand{\limb}{\lim\limits_}
\newcommand{\lebn}{\lambda_n}
\newcommand{\pp}{\text{~presque partout}}
\newcommand{\derivp}[2]{\frac{\partial #1}{\partial #2}}
\newcommand{\famiu}{(u_i)_{i \in I}}
\newcommand{\sev}{\text{~s.e.v.}}
\newcommand{\Mnp}{M_{n,p}(\K)}
\newcommand{\Mn}{M_{n}(\K)}
\newcommand{\tq}{\text{~t.q.~}}
\newcommand{\mq}{~montrer~que~}
\newcommand{\et}{\quad \text{et} \quad}
\newcommand{\ou}{\text{~ou~}}
\newcommand{\Kx}{\K[X]}


%Commandes de gars chelou
\renewcommand{\subset}{\subseteq}

%Commande restriction, ty duckduckgo
\def\restr#1#2{\mathchoice
              {\setbox1\hbox{${\displaystyle #1}_{\scriptstyle #2}$}
              \restraux{#1}{#2}}
              {\setbox1\hbox{${\textstyle #1}_{\scriptstyle #2}$}
              \restraux{#1}{#2}}
              {\setbox1\hbox{${\scriptstyle #1}_{\scriptscriptstyle #2}$}
              \restraux{#1}{#2}}
              {\setbox1\hbox{${\scriptscriptstyle #1}_{\scriptscriptstyle #2}$}
              \restraux{#1}{#2}}}
\def\restraux#1#2{{#1\,\smash{\vrule height .8\ht1 depth .85\dp1}}_{\,#2}} 

%Quelques normes
\DeclarePairedDelimiter\abs{\lvert}{\rvert}
\DeclarePairedDelimiter\labs{\bigl\lvert}{\bigr\rvert}
\DeclarePairedDelimiter\norm{\lVert}{\rVert}
\DeclarePairedDelimiterXPP\normi[1]{}\lVert\rVert{_1}{\ifblank{#1}{\:\cdot\:}{#1}}
\DeclarePairedDelimiterXPP\normd[1]{}\lVert\rVert{_2}{\ifblank{#1}{\:\cdot\:}{#1}}
\DeclarePairedDelimiterXPP\normp[1]{}\lVert\rVert{_p}{\ifblank{#1}{\:\cdot\:}{#1}}
\DeclarePairedDelimiterXPP\normq[1]{}\lVert\rVert{_q}{\ifblank{#1}{\:\cdot\:}{#1}}
\DeclarePairedDelimiterXPP\norminf[1]{}\lVert\rVert{_\infty}{\ifblank{#1}{\:\cdot\:}{#1}}

%Bracket en angle
\DeclarePairedDelimiter\angl{\langle}{\rangle}

%Set, parenthèses
\newcommand{\set}[1]{\{ #1 \}}
\newcommand{\bigset}[1]{\bigl\{ #1 \bigr\}}
\newcommand{\Bigset}[1]{\Bigl\{ #1 \Bigr\}}
\newcommand{\bigpar}[1]{\bigl( #1 \bigr)}
\newcommand{\Bigpar}[1]{\Bigl( #1 \Bigr)}

%Opitmisation des opérateurs de comparaison
\DeclarePairedDelimiter\tio{o (}{)}
\DeclarePairedDelimiter\gro{O (}{)}


\pagestyle{fancy}

\fancyhead[C]{Théorème d'Abel radial et taubérien faible}
\fancyhead[R]{}

\begin{document}

\underline{Source :}
\begin{itemize}
\item Rombaldi - Analyse - page 384 (Abel radial) 
\item Gourdon - Analyse - page 264 (Taubérien faible)
\end{itemize} 

\underline{Leçons :}
\begin{itemize}

\item 243 : Séries entières, propriétés de la somme. Exemples et applications.
\item 241 : Suites et séries de fonctions. Exemples et contre-exemples.

\end{itemize}

\begin{nota}
On considère une série entière $\sum a_n z^n$ à coefficients complexes, de rayon de convergence fini $R \in ]0, \infty[$ et dont 
la somme est notée $f(z)$ pour $z \in D(0, R)$.
\end{nota}

\begin{thm}[Abel radial]
On suppose qu'il existe $z_0 \in S(0,R)$ tel que la série $\sum a_n z_0^n$ soit convergente. 
Dans ce cas, on a :
\[ \lim_{t \to 1^{-}} \sum_{n=0}^{\infty} a_n z_0^n t^n = \sum_{n=0}^{\infty} a_n z_0^n . \]
\end{thm}

\begin{proof}
Pour tous $n \in \Ne, t \in [0,1]$ on note :
\[ S_n(t) = \sum_{k=0}^n a_k z_0^k t^k \et \varphi(t) = f(tz_0) = \sum_{n=0}^{\infty} a_n z_0^n t^n  . \]
En effectuant une transformation d'Abel, on trouve :
\begin{align*}
S_n(t) & = a_0 + \sum_{k=1}^n a_k z_0^k t^k \\
& = S_0(1) + \sum_{k=1}^n (S_k(1)-S_{k-1}(1)) t^k \\
& = S_0(1) t^0 + \sum_{k=1}^n S_k(1) t^k - \sum_{k=0}^{n-1} S_k t^{k+1} \\
& = \sum_{k=0}^{n-1} S_k(1)(t^k-t^{k+1}) + S_n(1)t^n \\
& = (1-t) \sum_{k=0}^{n-1} S_k(1) t^k + S_n(1)t^n \qquad (*)
\end{align*}
Si $t \in [0,1[$, on sait que $\lim_{n \to \infty} S_n(1)t^t = \lim_{n \to \infty} \varphi(1) t = 0$, 
donc, par passage à la limite dans (*), on en déduit que $\sum S_k(1) t^k$ est fini (pour $t <1$) et que :
\[ \varphi(t) = (1-t) \sum_{n=0}^{\infty} S_n(1) t^n \qquad (t \in [0,1[) . \]
En utilisant l'égalité $\sum_{\N} t^n = \dfrac1{1-t}$ pour $t \in [0,1[$, on peut écrire :
\begin{align*}
\varphi(t) - \varphi(1) & = (1-t) \sum_{n=0}^{\infty} S_n(1) t^n - \varphi(1) (1-t) \sum_{\N} t^n \\
& = (1-t) \sum_{n=0}^{\infty} (S_n(1) - \varphi(1)) t^n \qquad (t \in [0,1[) . 
\end{align*}
Comme $\lim_{n \to \infty} S_n(1) = \varphi(1)$, pour tout $\veps >0$, il existe un entier $n_{\veps}$ tel que :
\[ \forall n \geq n_{\veps}, \abs*{S_n(1) - \varphi(1)} < \veps , \]
ce qui donne :
\begin{align*}
\abs{\varphi(t) - \varphi(1)} & < (1-t) \left ( \sum_{k=0}^{n_{\veps}} \abs*{S_k(1) - \varphi(1)} t^k + 
\veps \sum_{k=n_{\veps +1}}^{\infty} t^k \right ) \\
& < (1-t) \sum_{k=0}^{n_{\veps}} \abs*{S_k(1) - \varphi(1)} + \veps (1-t) \sum_{k=0}^{\infty} t^k \\
& < (1-t) A_{\veps} + \veps  \qquad (t \in [0,1[),
\end{align*}
en posant $A_{\veps} = \sum_{k=0}^{n_{\veps}} \abs*{S_k(1) - \varphi(1)}$. 
En prenant $t$ dans $[\delta_{\veps},1[$ de façon à ce que $(1-t) \leq \dfrac{\veps}{A_{\veps}}$ on a :
\[ \abs{\varphi(t) - \varphi(1)} < 2\veps , \]
et ceci étant possible pour toute valeur de $\veps$, on en déduit que :
\[ \lim_{t \to 1^{-}} \varphi(t) = \varphi(1) , \]
ce qui est exactement le résultat énoncé.
\end{proof}

\begin{rmq}
Ce théorème est un résultat de prolongement par continuité : si une fonction $f$ développable en série entière sur un disque $D(0, R)$ 
est telle que son développement converge en un point du cercle de rayon $R$, alors la restriction de $f$ à $]-R, R[$ se prolonge par
continuité en $R$.
\end{rmq}

\begin{rmq}
A noter que la réciproque du théorème radial d'Abel est fausse, par exemple, on a :
\[ \lim_{t \to 1^{-1}} \sum_{n=0}^{\infty} (-1)^n t^n = \lim_{t \to 1^{-1}}  \dfrac1{1+t} = \dfrac12 , \]
or la série $\sum (-1)^n$ diverge grossièrement. Cependant, une réciproque partielle est donnée par le théorème taubérien faible.
\end{rmq}


\begin{thm}[Taubérien faible]
On suppose que $R=1$. Si :
\[ \lim_{t \to 1^{-}} \sum_{n=0}^{\infty} a_n t^n = S \in \Cx \et a_n = o_{\infty} \left(\dfrac1{n} \right) , \]
alors $\sum_{n=0}^{\infty} a_n = S$.
\end{thm}

\begin{proof}
L'hypothèse $a_n = o_{\infty} \left(\dfrac1{n} \right)$ permet d'affirmer que $\lim_{n \to \infty} n \abs{a_n} = 0$. 
Ainsi, la suite $(n \abs{a_n})_\N$ est majorée et on note $M$ son majorant.
On note $S_n = \sum_{k=0}^n a_n$ pour tout $n \in \N$. On a :
\[ \forall n \in \Ne, \forall x \in ]0,1[,  , \quad S_n - f(x) = \sum_{k=1}^n a_k (1-x^k) - \sum_{k=n+1}^{\infty} a_k x^k . \]
Or $1-x^k = (1-x)(1+x+\dots + x^{k-1}) \leq k(1-x)$ pour $x \in ]0,1[$, d'où on déduit :
\begin{align*}
\abs{S_n - f(x)} & \leq (1-x) \sum_{k=1}^n k \abs{a_k} + \sum_{k=n+1}^{\infty} \dfrac{k \abs{a_k}}{n} x^k \\
& \leq (1-x) Mn + \dfrac{\sup_{k>n} k \abs{a_k}}{n(1-x)} .
\end{align*}
Fixons $\veps \in ]0,1[$, en substituant $x=1-\veps/n$ dans l'inégalité précédente, on obtient :
\[ \forall n \in \Ne, \quad \abs*{S_n - f \left (1 - \dfrac{\veps}{n} \right ) } \leq M \veps + \dfrac{\sup_{k>n} k \abs{a_k}}{\veps} . \]
Prenons maintenant $N_0$ un rang à partir duquel on a $n \abs{a_n} < \veps^2$ (c'est possible car $n \abs{a_n} \to 0$), on en déduit :
\[ \forall n \geq N_0, \quad \abs*{S_n - f \left (1 - \dfrac{\veps}{n} \right ) } \leq M \veps + \veps = (M+1) \veps . \]
Or, par hypothèse, $f(x) \to S$ pour $x \to 1^{-1}$, on peut donc trouver un rang $N_1$ tel que :
\[ \forall n \geq N_1, \quad \abs*{S - f \left (1 - \dfrac{\veps}{n} \right )} < \veps , \]
donc, en imposant de plus $N_1 \geq N_0$, on déduit :
\begin{align*}
\abs{S_n - S} & \leq \abs*{S_n - f \left (1 - \dfrac{\veps}{n} \right )} + \abs*{S - f \left (1 - \dfrac{\veps}{n} \right )} \\
& < (M+1) \veps + \veps = (M+2) \veps ,
\end{align*}
et ceci étant vrai pour $\veps \in ]0, 1[$, on en déduit que $S_n \to S$, d'où le résultat.
\end{proof}

\end{document}